\chapter{Representation and Belief Targets}
\label{ch:representation}

\section{Design Goals}

The observation representation must satisfy four constraints.
\begin{enumerate}
  \item Fixed structure each turn so batching and attention masks are simple.
  \item Explicit separation of field state, player party, and opponent party.
  \item Compatibility with first-person partial observations from Metamon.
  \item Discrete tokens suitable for embedding lookup and vocabulary growth.
\end{enumerate}

\section{Thirteen-Token Turn Layout}

Each turn is encoded as exactly thirteen discrete tokens.

\begin{table}[ht]
  \centering
  \caption{Turn token layout. Indices are zero-based.}
  \label{tab:token-layout}
  \begin{tabular}{cl}
    \toprule
    Index & Role \\
    \midrule
    0 & Field (weather, terrain, hazards on both sides) \\
    1--6 & Player party slots (padded to six) \\
    7--12 & Opponent party slots (partial info when unrevealed) \\
    \bottomrule
  \end{tabular}
\end{table}

The field token concatenates sorted effect keys for weather, terrain, and side
conditions. Each party token encodes species, discretized HP bucket, status,
active flag, and fainted flag. Unfilled slots use a padding token.

A battle is a sequence of turns. Training samples use a causal context window of
$N$ prior turns plus the current turn. Left padding aligns the most recent turn
to the final sequence position.

\section{History Window}

Given turn frames $x_1, \ldots, x_t$, we form context
\[
  C_t = \mathrm{pad}\bigl(x_{t-N+1}, \ldots, x_t\bigr)
\]
with a binary turn mask marking valid versus padded positions. The model may
not attend to future turns. This implements a finite-history approximation to
belief updating in a POMDP.

\section{Belief Supervision Targets}

Belief heads predict hidden opponent attributes knowable from full logs but not
from the first-person view at time $t$. Primary targets in priority order are
\begin{enumerate}
  \item species identity for each unrevealed opponent slot,
  \item moves for revealed or partially revealed opponent Pok\'emon,
  \item held item when inferable before reveal,
  \item ability when not yet shown,
  \item terastallization type in Gen 9 OU.
\end{enumerate}

Each target is a multi-class prediction per opponent slot. Revealed or empty
slots are masked out of the loss. Labels are extracted from Metamon full-state
annotations at each turn boundary.

For species belief we define a vocabulary over base species names. Slot $s
\in \{0,\ldots,5\}$ maps to opponent token index $7+s$ and receives a label
$y_{t,s} \in \mathcal{Y}_\mathrm{species} \cup \{\emptyset\}$ where $\emptyset$
denotes masked slots.

\section{Combinatorial Keys for Evaluation}

We define two keys for distribution shift experiments.
\begin{description}
  \item[Team key] Sorted tuple of six species names for the opponent roster.
  \item[Set key] Tuple of species, item, ability, and move quartet for a slot.
\end{description}

Held-out keys define test battles where the opponent configuration never
appeared in training. This tests compositional generalization beyond random
battle splits.

\section{Implementation}

The \texttt{TurnTokenizer} class in \texttt{src/pokemon\_thesis/tokenizer/}
implements encoding from JSON battle records. Encoded trajectories are stored as
JSONL with \texttt{token\_ids}, \texttt{action}, and optional \texttt{belief}
fields for multi-task training.
